\documentclass[10pt,a4paper]{amsart}
\usepackage{amsmath,amssymb,amsthm,graphicx}
\usepackage{hyperref}
% Minimal class-compatibility preamble for a reader excerpt of
%   Bhimarthi Ravinder, "Demazure Modules, Chari-Venkatesh Modules and Fusion Products",
%   SIGMA 10 (2014), 110, DOI 10.3842/SIGMA.2014.110.
%
% The author's pinned TeX file declares no theorem environments of its own: they are
% supplied by the SIGMA publisher class sigma.cls, which is not installed here.  The
% block below is that class's theorem-environment block (same environment names and
% styles), copied mechanically so that the retained author text compiles unchanged.
% No mathematical macro or statement is introduced by this file.
\newtheorem{theorem}{Theorem}
\newtheorem*{theorem*}{Theorem}
\newtheorem{lemma}{Lemma}
\newtheorem*{lemma*}{Lemma}
\newtheorem{corollary}{Corollary}
\newtheorem*{corollary*}{Corollary}
\newtheorem{proposition}{Proposition}
\newtheorem*{proposition*}{Proposition}
\newtheorem{conjecture}{Conjecture}
\newtheorem*{conjecture*}{Conjecture}
{\theoremstyle{definition} \newtheorem{definition}{Definition}
\newtheorem*{definition*}{Definition}
\newtheorem{example}{Example}
\newtheorem*{example*}{Example}
\newtheorem{remark}{Remark}
\newtheorem*{remark*}{Remark}
\newtheorem{note}{Note}
\newtheorem*{note*}{Note}
}
\DeclareMathOperator{\ev}{ev}
\DeclareMathOperator{\gr}{gr}
\DeclareMathOperator{\spn}{span}
\title{Relations satisfied by the shifted highest-root vector in the level-one Demazure module D(1,(k+1)theta)}
\author{Bhimarthi Ravinder}
\date{doi:10.3842/SIGMA.2014.110; CC BY 4.0}
\begin{document}
\maketitle
\noindent\textit{Editorial note.} All mathematical text is the author's own, copied from the fixed publisher source; the proof appears exactly as printed and nothing was rewritten.\par\medskip
\noindent\textit{Editorial note.} The author's file declares no theorem environments of its own: they come from the SIGMA publisher class, which is not installed for this excerpt. The minimal preamble therefore reproduces that class's theorem-environment block (same environment names and styles); the author's own operator macros are taken from source lines 9-11.\par\medskip
\noindent\textit{Editorial note.} Context blocks are the source material the proof uses: Garland's quoted identity, source lemma g, the definition of the level-one Demazure module, and the preceding source proposition WvsD (whose defining relations for D(1,(k+1)theta) are used at lines 489-490). They are marked as context and are not counted as this task's proof. Environment numbers printed here are local to the excerpt; the source statement is the lemma labelled crucial at source line 454 (Section 3).\par\medskip
\noindent\textit{Allowed standard prerequisites (editorial).}\begin{itemize}
\item Root systems and sl\_2-triple arguments in a finite-dimensional complex simple Lie algebra. Roots, coroots, the highest root theta and d\_alpha = 2/(alpha|alpha) are assumed. The weight computation for (x\textasciicircum{}-\_theta (x) t\textasciicircum{}\{2k+1-i\}) wbar, the vanishing of its x\textasciicircum{}+\_alpha (x) 1 image via theta - alpha in R\textasciicircum{}+, and the resulting sl\_2 conclusion are retained obligations.
\item Garland's current-algebra identity, as quoted by the author. The identity is quoted verbatim in the excerpt and is not re-proved. Its specialisation r = 2, s = 3k+1-i, alpha = theta and the subsequent manipulation of the resulting sum are the retained proof.
\item Defining relations of the level-one Demazure module, obtained from the preceding source proposition. The excerpt carries the definition of D(1,lambda) and the author's presentation of D(1,k theta). The proof uses the resulting relations of D(1,(k+1)theta) at lines 489-490; nothing else about Demazure modules is assumed.
\end{itemize}\medskip
\begin{lemma}\label{garland}
Given $s\in\mathbb{N}$, $r\in\mathbb{Z}_{\geq0}$ and $\alpha\in R^{+}$, we have
\begin{gather*}
(x^{+}_{\alpha} \otimes t)^{(s)}(x^{-}_{\alpha} \otimes
1)^{(s+r)}-(-1)^s\mathbf{x}^{-}_{\alpha}(r,s)\in\mathbf{U}(\mathfrak{g}[t])\mathfrak{n}^{+}[t]\bigoplus
\mathbf{U}(\mathfrak{n}^{-}[t])\mathfrak{h}[t]_{+}.
\end{gather*}
\end{lemma}

We now specialize to the case $\lambda\in \mathbb{N}\theta$, and obtain some further useful relations that hold in~$W(\lambda)$.
\begin{lemma}
\label{g}
Let $k\in\mathbb{N}$.
The following relations hold in the local Weyl module $W((k+1)\theta)$:
\begin{enumerate}\itemsep=0pt
\item[$1)$] $(x^{-}_{\theta}\otimes 1)^{2k+1}\big(x^{-}_{\theta}\otimes t^{2k+1-i}\big)w_{(k+1)\theta}=0,
\qquad
\forall\,
0\leq i \leq k$;
\item[$2)$] $(x^{-}_{\theta}\otimes t^m)(x^{-}_{\theta} \otimes t^{m+1})  w_{(k+1)\theta} \in  \big\langle \big(x^{-}_{\theta}
\otimes t^{m+2}\big) w_{(k+1)\theta} \rangle,
\qquad
\forall\,
m\geq k$.
\end{enumerate}
\end{lemma}

\begin{definition}
\textup{(see~\cite[Corollary 3.5]{CV})} The level one Demazure module $D(1,\lambda)$ is the graded quotient of
$W(\lambda)$ by the submodule generated by the union of the following two sets:
\begin{gather}
\big\{(x^{-}_{\alpha} \otimes t^{s_{\alpha}})  w_{\lambda}: \alpha\in R^{+}~\textup{such that}~d_{\alpha} > 1\big\},
\label{dm1}
\\
\big\{\big(x^{-}_{\alpha} \otimes t^{s_{\alpha}-1}\big)^2  w_{\lambda}: \alpha\in R^{+}~\textup{such that}~d_\alpha =3~\textup{and}~m_\alpha=1\big\}.
\label{dm2}
\end{gather}
In particular, for $\mathfrak{g}$ simply laced, $D(1,\lambda)\cong_{\mathfrak{g}[t]} W(\lambda)$. We denote~by
$\overline{w}_\lambda$, the image of $w_\lambda$ in~$D(1,\lambda)$.
\end{definition}

The following proposition gives explicit def\/ining relations for $D(1,k\theta)$.
\begin{proposition}
\label{WvsD}
Given $k\geq 1$, the level~$1$ Demazure module $D(1,k\theta)$ is the graded $\mathfrak{g}[t]$-module generated by an
element $\overline{w}_{k\theta}$, with the following defining relations:
\begin{gather*}
\mathfrak{n}^{+}[t]\, \overline{w}_{k\theta}=0,
\qquad
(h \otimes t^s) \overline{w}_{k\theta} = \langle k\theta, h\rangle \delta_{s,0} \overline{w}_{k\theta},
\qquad
s\geq0,
\quad
h\in\mathfrak{h},
\\
(x^{-}_{\alpha}\otimes 1) \overline{w}_{k\theta}=0,
\qquad
\alpha\in R^+,
\quad
(\theta|\alpha)=0,
\\
(x^{-}_{\alpha} \otimes 1)^{kd_{\alpha}+1} \overline{w}_{k\theta}=0,
\qquad
\big(x^{-}_{\alpha} \otimes t^k\big) \overline{w}_{k\theta}=0,
\qquad
\alpha\in R^+,
\quad
(\theta|\alpha)=1,
\\
(x^{-}_{\theta} \otimes 1)^{2k+1} \overline{w}_{k\theta}=0.
\end{gather*}
\end{proposition}

\begin{lemma}
\label{crucial}
Let $k\geq 1$ and $0\leq i \leq k$. The following relations hold in the module $D(1,(k+1)\theta)$:
\begin{enumerate}\itemsep=0pt
\item[$1)$] $(x^{-}_{\alpha} \otimes 1)^{kd_{\alpha}+1}\big(x^{-}_{\theta} \otimes t^{2k+1-i}\big) \overline{w}_{(k+1)\theta}=0$,
$\forall\, \alpha \in R^{+}$, $(\theta|\alpha)=1$;
\item[$2)$] $\big(x^{-}_{\alpha} \otimes t^{k}\big) (x^{-}_{\theta} \otimes t^{2k+1-i}) \overline{w}_{(k+1)\theta} \in
\big\langle \big(x^{-}_{\theta} \otimes t^{2k+2-i}\big) \overline{w}_{(k+1)\theta} \big\rangle$,
$\forall\, \alpha \in R^{+}$, $(\theta|\alpha)=1$;
\item[$3)$] $\big(x^{-}_{\theta} \otimes t^{2k-i}\big) \big(x^{-}_{\theta} \otimes t^{2k+1-i}\big) \overline{w}_{(k+1)\theta}
\in
\big\langle \big(x^{-}_{\theta} \otimes t^{2k+2-i}\big) \overline{w}_{(k+1)\theta} \big\rangle$.
\end{enumerate}
\end{lemma}

\begin{proof}
Let $ \alpha \in R^{+} \textup{with} (\theta|\alpha)=1$. This implies that $\theta-\alpha$ is also a~root of
$\mathfrak{g}$ and $(\theta|\theta-\alpha)=1$. We now prove part (1).
Observe that, $(x^{-}_{\theta} \otimes t^{2k+1-i}) \overline{w}_{(k+1)\theta}$ is an element of weight $k\theta$.
Further $(x^{+}_{\alpha} \otimes 1) (x^{-}_{\theta} \otimes t^{2k+1-i}) \overline{w}_{(k+1)\theta}=0$, since
$(x^{+}_{\alpha} \otimes 1)\overline{w}_{(k+1)\theta}=0$ and $(x^{-}_{\theta-\alpha} \otimes t^{2k+1-i}) \overline{w}_{(k+1)\theta}=0$,
for all $0\leq i \leq k$.
Considering the copy of $\mathfrak{sl}_2$ spanned by $x^{+}_{\alpha} \otimes 1$, $x^{-}_{\alpha} \otimes 1$,
$\alpha^{\vee}\otimes 1$, we obtain part (1) by standard $\mathfrak{sl}_2$ arguments.
We now prove part (2).
Putting $r=2$, $s=(3k+1-i)$ and $\alpha=\theta$ in Lemma~\ref{garland}, we get
\begin{gather}
\big(x^{-}_{\theta} \otimes t^{k}\big) \big(x^{-}_{\theta} \otimes t^{2k+1-i}\big) \overline{w}_{(k+1)\theta} +
\sum\limits_{\substack{k+1\leq p\leq q \leq 2k-i\\p+q= 3k+1-i}}
\frac{1}{(2\delta_{p, q})!} \big(x^{-}_{\theta} \otimes t^{p}\big) \big(x^{-}_{\theta} \otimes t^{q}\big)
\overline{w}_{(k+1)\theta}
\nonumber\\
\qquad{}
\in
\big\langle\big(x^{-}_{\theta}\otimes t^{2k+2-i}\big) \overline{w}_{(k+1)\theta} \big\rangle, \label{e}
\end{gather}
since $(x^{-}_{\theta}\otimes 1)^{3k+3-i}\overline{w}_{(k+1)\theta}=0$, $\forall\, 0\leq i \leq k$.
Now we act on both sides of~\eqref{e} by $x^{+}_{\theta-\alpha}$ and use the relation $(x^{-}_{\alpha} \otimes t^r) \overline{w}_{(k+1)\theta}=0$,
for all $r\geq(k+1)$, which gives part~(2).
Part~(3) is immediate from the part~(2) of Lemma~\ref{g}.
\end{proof}

% Bibliography of the delivered excerpt: the author's own thebibliography,
% copied verbatim from pinned-original-source.tex lines 875-924.
% The SIGMA-class-only end-of-article token (line 924) is omitted because the
% excerpt is typeset with a class-compatibility preamble instead of sigma.cls.
\begin{thebibliography}{99}
\footnotesize \itemsep=0pt

\bibitem{CFK}
Chari V., Fourier G., Khandai T., A categorical approach to {W}eyl modules,
  \href{http://dx.doi.org/10.1007/s00031-010-9090-9}{\textit{Transform. Groups}} \textbf{15} (2010), 517--549, \href{http://arxiv.org/abs/0906.2014}{arXiv:0906.2014}.

\bibitem{CP}
Chari V., Pressley A., Weyl modules for classical and quantum af\/fine algebras,
  \href{http://dx.doi.org/10.1090/S1088-4165-01-00115-7}{\textit{Represent. Theory}} \textbf{5} (2001), 191--223,
  \href{http://arxiv.org/abs/math.QA/0004174}{math.QA/0004174}.

\bibitem{CSVW}
Chari V., Shereen P., Venkatesh R., Wand J., A Steinberg type decomposition
  theorem for higher level Demazure modules, \href{http://arxiv.org/abs/1408.4090}{arXiv:1408.4090}.

\bibitem{CV}
Chari V., Venkatesh R., Demazure modules, fusion products, and $Q$-systems,
  \href{http://dx.doi.org/10.1007/s00220-014-2175-x}{\textit{Comm. Math. Phys.}}, {t}o appear, \href{http://arxiv.org/abs/1305.2523}{arXiv:1305.2523}.

\bibitem{FL}
Feigin B., Loktev S., On generalized {K}ostka polynomials and the quantum
  {V}erlinde rule, in Dif\/ferential Topology, Inf\/inite-Dimensional {L}ie
  Algebras, and Applications, \textit{Amer. Math. Soc. Transl. Ser.~2}, Vol.~194, Amer. Math. Soc., Providence, RI, 1999, 61--79,
  \href{http://arxiv.org/abs/math.QA/9812093}{math.QA/9812093}.

\bibitem{F}
Feigin E., The {PBW} f\/iltration, {D}emazure modules and toroidal current
  algebras, \href{http://dx.doi.org/10.3842/SIGMA.2008.070}{\textit{SIGMA}} \textbf{4} (2008), 070, 21~pages,
  \href{http://arxiv.org/abs/0806.4851}{arXiv:0806.4851}.

\bibitem{FoL}
Fourier G., Littelmann P., Weyl modules, {D}emazure modules, {KR}-modules,
  crystals, fusion products and limit constructions, \href{http://dx.doi.org/10.1016/j.aim.2006.09.002}{\textit{Adv. Math.}}
  \textbf{211} (2007), 566--593, \href{http://arxiv.org/abs/math.RT/0509276}{math.RT/0509276}.

\bibitem{G}
Garland H., The arithmetic theory of loop algebras, \href{http://dx.doi.org/10.1016/0021-8693(78)90294-6}{\textit{J.~Algebra}}
  \textbf{53} (1978), 480--551.

\bibitem{Naoi}
Naoi K., Weyl modules, {D}emazure modules and f\/inite crystals for non-simply
  laced type, \href{http://dx.doi.org/10.1016/j.aim.2011.10.005}{\textit{Adv. Math.}} \textbf{229} (2012), 875--934,
  \href{http://arxiv.org/abs/1012.5480}{arXiv:1012.5480}.

\bibitem{V}
Venkatesh R., Fusion product structure of Demazure modules, \href{http://dx.doi.org/10.1007/s10468-014-9495-6}{\textit{Algebr.
  Represent. Theory}}, {t}o appear, \href{http://arxiv.org/abs/1311.2224}{arXiv:1311.2224}.

\end{thebibliography}
\end{document}
